% ============================================================
% Question Bank: ch02-10 Compound Interest Introduction
% Topic Title: Compound Interest Introduction
% CCSS: Supplementary
% Questions: 30 (18 MC + 7 SA + 5 Graphical)
% ============================================================

% --- Multiple Choice Questions (18) ---

\begin{practiceQuestion}{2-10-q01}{mc}
\begin{questionText}
What is the main difference between simple and compound interest?
\end{questionText}
\choiceA{Simple interest uses a higher rate}
\choiceB{Compound interest is calculated on the principal plus previously earned interest}
\choiceC{Simple interest compounds monthly}
\choiceD{Compound interest only applies to loans}
\correctAnswer{B}
\explanation{Compound interest is calculated on the principal plus any interest already earned, so your interest earns interest.}
\end{practiceQuestion}

\begin{practiceQuestion}{2-10-q02}{mc}
\begin{questionText}
You deposit $\$1{,}000$ at $10\%$ compounded yearly. What is the total after $1$ year?
\end{questionText}
\choiceA{$\$1{,}010$}
\choiceB{$\$1{,}100$}
\choiceC{$\$1{,}200$}
\choiceD{$\$1{,}001$}
\correctAnswer{B}
\explanation{$1{,}000 \times 1.10 = \$1{,}100$.}
\end{practiceQuestion}

\begin{practiceQuestion}{2-10-q03}{mc}
\begin{questionText}
You deposit $\$500$ at $10\%$ compounded yearly. What is the total after $2$ years?
\end{questionText}
\choiceA{$\$600$}
\choiceB{$\$605$}
\choiceC{$\$550$}
\choiceD{$\$610$}
\correctAnswer{B}
\explanation{Year 1: $500 \times 1.10 = \$550$. Year 2: $550 \times 1.10 = \$605$.}
\end{practiceQuestion}

\begin{practiceQuestion}{2-10-q04}{mc}
\begin{questionText}
$\$400$ is deposited at $5\%$ simple interest for $3$ years. How much interest is earned?
\end{questionText}
\choiceA{$\$20$}
\choiceB{$\$60$}
\choiceC{$\$63.05$}
\choiceD{$\$15$}
\correctAnswer{B}
\explanation{Simple interest: $I = 400 \times 0.05 \times 3 = \$60$.}
\end{practiceQuestion}

\begin{practiceQuestion}{2-10-q05}{mc}
\begin{questionText}
$\$400$ is deposited at $5\%$ compounded yearly for $3$ years. What is the total amount?
\end{questionText}
\choiceA{$\$460$}
\choiceB{$\$460.60$}
\choiceC{$\$463.05$}
\choiceD{$\$475$}
\correctAnswer{C}
\explanation{Use $A = P(1+r)^t = 400(1.05)^3 = 400 \times 1.157625 = \$463.05$.}
\end{practiceQuestion}

\begin{practiceQuestion}{2-10-q06}{mc}
\begin{questionText}
In the formula $A = P(1 + r)^t$, what does $P$ represent?
\end{questionText}
\choiceA{Percent}
\choiceB{Principal (starting amount)}
\choiceC{Payment}
\choiceD{Profit}
\correctAnswer{B}
\explanation{$P$ is the principal --- the starting amount of money deposited or borrowed.}
\end{practiceQuestion}

\begin{practiceQuestion}{2-10-q07}{mc}
\begin{questionText}
You deposit $\$200$ at $8\%$ compounded yearly. What is the amount after $1$ year?
\end{questionText}
\choiceA{$\$208$}
\choiceB{$\$216$}
\choiceC{$\$280$}
\choiceD{$\$216.32$}
\correctAnswer{B}
\explanation{$200 \times 1.08 = \$216$.}
\end{practiceQuestion}

\begin{practiceQuestion}{2-10-q08}{mc}
\begin{questionText}
Deposit $\$600$ at $10\%$ for $2$ years. How much MORE does compound interest earn compared to simple interest?
\end{questionText}
\choiceA{$\$6$}
\choiceB{$\$12$}
\choiceC{$\$120$}
\choiceD{$\$60$}
\correctAnswer{A}
\explanation{Simple: $I = 600 \times 0.10 \times 2 = \$120$, total $= \$720$. Compound: Year 1: $660$; Year 2: $660 \times 1.10 = \$726$. Difference: $726 - 720 = \$6$.}
\end{practiceQuestion}

\begin{practiceQuestion}{2-10-q09}{mc}
\begin{questionText}
Which grows faster over time?
\end{questionText}
\choiceA{Simple interest}
\choiceB{Compound interest}
\choiceC{They grow at the same rate}
\choiceD{It depends on the principal}
\correctAnswer{B}
\explanation{Compound interest grows faster because you earn interest on your interest, making the base larger each year.}
\end{practiceQuestion}

\begin{practiceQuestion}{2-10-q10}{mc}
\begin{questionText}
You invest $\$1{,}000$ at $6\%$ compounded yearly. After $2$ years, the total is:
\end{questionText}
\choiceA{$\$1{,}120$}
\choiceB{$\$1{,}123.60$}
\choiceC{$\$1{,}060$}
\choiceD{$\$1{,}126$}
\correctAnswer{B}
\explanation{Year 1: $1{,}000 \times 1.06 = \$1{,}060$. Year 2: $1{,}060 \times 1.06 = \$1{,}123.60$.}
\end{practiceQuestion}

\begin{practiceQuestion}{2-10-q11}{mc}
\begin{questionText}
$\$300$ at $10\%$ simple interest for $3$ years. What is the total?
\end{questionText}
\choiceA{$\$330$}
\choiceB{$\$390$}
\choiceC{$\$399.30$}
\choiceD{$\$360$}
\correctAnswer{B}
\explanation{$I = 300 \times 0.10 \times 3 = \$90$. Total $= 300 + 90 = \$390$.}
\end{practiceQuestion}

\begin{practiceQuestion}{2-10-q12}{mc}
\begin{questionText}
$\$300$ at $10\%$ compounded yearly for $3$ years. What is the total?
\end{questionText}
\choiceA{$\$390$}
\choiceB{$\$393$}
\choiceC{$\$399.30$}
\choiceD{$\$400$}
\correctAnswer{C}
\explanation{Year 1: $300 \times 1.10 = 330$. Year 2: $330 \times 1.10 = 363$. Year 3: $363 \times 1.10 = \$399.30$.}
\end{practiceQuestion}

\begin{practiceQuestion}{2-10-q13}{mc}
\begin{questionText}
In the formula $A = P(1 + r)^t$, what does $t$ represent?
\end{questionText}
\choiceA{Tax rate}
\choiceB{Total amount}
\choiceC{Number of years}
\choiceD{Type of account}
\correctAnswer{C}
\explanation{$t$ is the number of years the money is invested or borrowed.}
\end{practiceQuestion}

\begin{practiceQuestion}{2-10-q14}{mc}
\begin{questionText}
A $\$700$ loan at $5\%$ compounded yearly for $2$ years. How much is owed?
\end{questionText}
\choiceA{$\$770$}
\choiceB{$\$771.75$}
\choiceC{$\$735$}
\choiceD{$\$772$}
\correctAnswer{B}
\explanation{Year 1: $700 \times 1.05 = \$735$. Year 2: $735 \times 1.05 = \$771.75$.}
\end{practiceQuestion}

\begin{practiceQuestion}{2-10-q15}{mc}
\begin{questionText}
Why is compound interest sometimes called ``interest on interest''?
\end{questionText}
\choiceA{The bank charges interest twice}
\choiceB{Each year the interest is calculated on the total, including past interest}
\choiceC{You earn interest only if you don't withdraw}
\choiceD{The rate doubles each year}
\correctAnswer{B}
\explanation{With compound interest, each year's interest is added to the total, so the next year's interest is larger.}
\end{practiceQuestion}

\begin{practiceQuestion}{2-10-q16}{mc}
\begin{questionText}
You deposit $\$250$ at $4\%$ compounded yearly. What is the balance after $2$ years?
\end{questionText}
\choiceA{$\$270$}
\choiceB{$\$270.40$}
\choiceC{$\$260$}
\choiceD{$\$280$}
\correctAnswer{B}
\explanation{Year 1: $250 \times 1.04 = \$260$. Year 2: $260 \times 1.04 = \$270.40$.}
\end{practiceQuestion}

\begin{practiceQuestion}{2-10-q17}{mc}
\begin{questionText}
Using $A = P(1+r)^t$, find the value of $\$1{,}000$ at $3\%$ for $2$ years.
\end{questionText}
\choiceA{$\$1{,}060$}
\choiceB{$\$1{,}060.90$}
\choiceC{$\$1{,}030$}
\choiceD{$\$1{,}090$}
\correctAnswer{B}
\explanation{$A = 1{,}000(1.03)^2 = 1{,}000 \times 1.0609 = \$1{,}060.90$.}
\end{practiceQuestion}

\begin{practiceQuestion}{2-10-q18}{mc}
\begin{questionText}
A $\$800$ deposit earns $\$48$ in simple interest over $2$ years. What is the annual interest rate?
\end{questionText}
\choiceA{$2\%$}
\choiceB{$3\%$}
\choiceC{$4\%$}
\choiceD{$6\%$}
\correctAnswer{B}
\explanation{$48 = 800 \times r \times 2$. So $r = \frac{48}{1{,}600} = 0.03 = 3\%$.}
\end{practiceQuestion}

% --- Short Answer Questions (7) ---

\begin{practiceQuestion}{2-10-q19}{sa}
\begin{questionText}
Deposit $\$600$ at $5\%$ compounded yearly. Find the total after $2$ years.
\end{questionText}
\correctAnswer{$\$661.50$}
\explanation{Year 1: $600 \times 1.05 = \$630$. Year 2: $630 \times 1.05 = \$661.50$.}
\end{practiceQuestion}

\begin{practiceQuestion}{2-10-q20}{sa}
\begin{questionText}
$\$900$ at $10\%$ simple interest for $3$ years. What is the total interest earned?
\end{questionText}
\correctAnswer{$\$270$}
\explanation{$I = 900 \times 0.10 \times 3 = \$270$.}
\end{practiceQuestion}

\begin{practiceQuestion}{2-10-q21}{sa}
\begin{questionText}
$\$900$ at $10\%$ compounded yearly for $3$ years. What is the total amount?
\end{questionText}
\correctAnswer{$\$1{,}197.90$}
\explanation{Year 1: $900 \times 1.10 = 990$. Year 2: $990 \times 1.10 = 1{,}089$. Year 3: $1{,}089 \times 1.10 = \$1{,}197.90$.}
\end{practiceQuestion}

\begin{practiceQuestion}{2-10-q22}{sa}
\begin{questionText}
How much more does compound interest earn than simple interest on $\$500$ at $8\%$ for $3$ years?
\end{questionText}
\correctAnswer{$\$9.86$}
\explanation{Simple: $I = 500 \times 0.08 \times 3 = \$120$, total $= \$620$. Compound: $A = 500(1.08)^3 = 500 \times 1.259712 = \$629.86$. Difference: $629.86 - 620 = \$9.86$.}
\end{practiceQuestion}

\begin{practiceQuestion}{2-10-q23}{sa}
\begin{questionText}
Use $A = P(1 + r)^t$ to find the value of $\$1{,}500$ at $6\%$ compounded yearly for $2$ years.
\end{questionText}
\correctAnswer{$\$1{,}685.40$}
\explanation{$A = 1{,}500(1.06)^2 = 1{,}500 \times 1.1236 = \$1{,}685.40$.}
\end{practiceQuestion}

\begin{practiceQuestion}{2-10-q24}{sa}
\begin{questionText}
A $\$400$ loan at $12\%$ compounds yearly for $2$ years. How much is owed?
\end{questionText}
\correctAnswer{$\$501.76$}
\explanation{Year 1: $400 \times 1.12 = \$448$. Year 2: $448 \times 1.12 = \$501.76$.}
\end{practiceQuestion}

\begin{practiceQuestion}{2-10-q25}{sa}
\begin{questionText}
Explain in your own words why compound interest grows faster than simple interest.
\end{questionText}
\correctAnswer{You earn interest on your previously earned interest}
\explanation{Each year, the interest earned is added to the principal. The next year, you earn interest on a larger amount, so it grows faster.}
\end{practiceQuestion}

% --- Graphical Questions (3 GMC + 2 GSA) ---

\begin{practiceQuestion}{2-10-q26}{gmc}
\begin{questionText}
The table shows the balance of a $\$1{,}000$ deposit at $10\%$ each year. Is this simple or compound interest?

\begin{center}
\begin{tabular}{|c|c|}
\hline
\textbf{Year} & \textbf{Balance} \\
\hline
$0$ & $\$1{,}000$ \\
\hline
$1$ & $\$1{,}100$ \\
\hline
$2$ & $\$1{,}210$ \\
\hline
$3$ & $\$1{,}331$ \\
\hline
\end{tabular}
\end{center}
\end{questionText}
\choiceA{Simple interest}
\choiceB{Compound interest}
\choiceC{Neither}
\choiceD{Cannot be determined}
\correctAnswer{B}
\explanation{Simple interest would add $\$100$ each year (totals: $1{,}100$; $1{,}200$; $1{,}300$). Here the amounts grow more each year, which means it is compound interest.}
\end{practiceQuestion}

\begin{practiceQuestion}{2-10-q27}{gmc}
\begin{questionText}
The line graph below shows two accounts over $3$ years, both starting at $\$500$ with $10\%$ interest. Which line represents compound interest?

\begin{center}
\begin{tikzpicture}[scale=0.8]
  \draw[thick,->] (0,0) -- (5,0) node[right] {Year};
  \draw[thick,->] (0,0) -- (0,5) node[above] {\$};
  \foreach \x in {0,1,2,3} {
    \draw (\x+0.5,-0.1) -- (\x+0.5,0.1);
    \node[below] at (\x+0.5,-0.15) {\footnotesize \x};
  }
  \foreach \y/\label in {1/500, 2/550, 3/600, 4/650} {
    \draw (-0.1,\y) -- (0.1,\y);
    \node[left] at (-0.15,\y) {\footnotesize \$\label};
  }
  % Simple (Line A) - straight line
  \draw[thick,funRed,mark=square*] plot coordinates {(0.5,1) (1.5,2) (2.5,3) (3.5,4)};
  \node[right,funRed] at (3.7,4) {\footnotesize Line A};
  % Compound (Line B) - curved
  \draw[thick,funBlue,mark=*] plot coordinates {(0.5,1) (1.5,2.1) (2.5,3.32) (3.5,4.65)};
  \node[right,funBlue] at (3.7,4.65) {\footnotesize Line B};
\end{tikzpicture}
\end{center}
\end{questionText}
\choiceA{Line A}
\choiceB{Line B}
\choiceC{Both lines}
\choiceD{Neither line}
\correctAnswer{B}
\explanation{Compound interest grows by increasing amounts each year, so the curve bends upward. Line B is compound; Line A (straight) is simple.}
\end{practiceQuestion}

\begin{practiceQuestion}{2-10-q28}{gmc}
\begin{questionText}
The table compares simple and compound interest on $\$200$ at $10\%$ for $3$ years. What is the compound interest total after Year $3$?

\begin{center}
\begin{tabular}{|c|c|c|}
\hline
\textbf{Year} & \textbf{Simple} & \textbf{Compound} \\
\hline
$0$ & $\$200$ & $\$200$ \\
\hline
$1$ & $\$220$ & $\$220$ \\
\hline
$2$ & $\$240$ & $\$242$ \\
\hline
$3$ & $\$260$ & ? \\
\hline
\end{tabular}
\end{center}
\end{questionText}
\choiceA{$\$260$}
\choiceB{$\$264$}
\choiceC{$\$266.20$}
\choiceD{$\$268$}
\correctAnswer{C}
\explanation{Year 3 compound: $242 \times 1.10 = \$266.20$.}
\end{practiceQuestion}

\begin{practiceQuestion}{2-10-q29}{gsa}
\begin{questionText}
The bar graph shows the total amount after $2$ years for a $\$1{,}000$ deposit. One bar shows simple interest; the other shows compound interest, both at $5\%$. Find the difference between the two totals.

\begin{center}
\begin{tikzpicture}[scale=0.7]
  \draw[thick,->] (0,0) -- (0,5.5) node[above] {\$};
  \draw[thick,->] (0,0) -- (6.5,0);
  \foreach \y/\label in {0/0, 1/200, 2/400, 3/600, 4/800, 5/1000} {
    \draw (-0.1,\y) -- (0.1,\y);
    \node[left] at (-0.15,\y) {\footnotesize \$\label};
  }
  % Simple bar
  \fill[funRed!70] (1,0) rectangle (2.5,5.5);
  \node[below] at (1.75,0) {\footnotesize Simple};
  \node[above] at (1.75,5.5) {\footnotesize \$1{,}100};
  % Compound bar
  \fill[funBlue!70] (3.5,0) rectangle (5,5.5125);
  \node[below] at (4.25,0) {\footnotesize Compound};
  \node[above] at (4.25,5.5125) {\footnotesize \$1{,}102.50};
\end{tikzpicture}
\end{center}
\end{questionText}
\correctAnswer{$\$2.50$}
\explanation{Simple: $1{,}000 + (1{,}000 \times 0.05 \times 2) = \$1{,}100$. Compound: $1{,}000 \times (1.05)^2 = \$1{,}102.50$. Difference: $1{,}102.50 - 1{,}100 = \$2.50$.}
\end{practiceQuestion}

\begin{practiceQuestion}{2-10-q30}{gsa}
\begin{questionText}
Complete the table for a $\$800$ deposit at $5\%$ compounded yearly. What is the balance at the end of Year $3$?

\begin{center}
\begin{tabular}{|c|c|c|}
\hline
\textbf{Year} & \textbf{Interest Earned} & \textbf{Balance} \\
\hline
$0$ & --- & $\$800$ \\
\hline
$1$ & $\$40$ & $\$840$ \\
\hline
$2$ & $\$42$ & $\$882$ \\
\hline
$3$ & ? & ? \\
\hline
\end{tabular}
\end{center}
\end{questionText}
\correctAnswer{$\$926.10$}
\explanation{Year 3 interest: $882 \times 0.05 = \$44.10$. Balance: $882 + 44.10 = \$926.10$.}
\end{practiceQuestion}
